% !PRINTED 681
\weberpage{681}
%<<CONT

2. \gesperrt{Seien $\alpha, \beta, \gamma$ drei Funktionen in $\Omega$; dann ist}
\weq{6}{\left(\frac{d\alpha}{d\beta}\right)\left(\frac{d\beta}{d\gamma}\right) = \left(\frac{d\alpha}{d\gamma}\right).}

Denn es ist für unendlich viele Punkte
\[
\left(\frac{\alpha-\alpha_0}{\beta-\beta_0}\right)_0 \left(\frac{\beta-\beta_0}{\gamma-\gamma_0}\right)_0 = \left(\frac{\alpha-\alpha_0}{\gamma-\gamma_0}\right)_0,
\]
und daraus folgt wie bei 1. der Satz 2.

Demnach kann man jeder Funktion $\alpha$ in $\Omega$ eine Funktion $d\alpha$ so zuordnen, daß für irgend zwei dieser Funktionen:
\weq{7}{\left(\frac{d\alpha}{d\beta}\right) = \frac{d\alpha}{d\beta}}
ein wirklicher Quotient zweier Funktionen $d\alpha, d\beta$ wird. Diese Funktionen nennen wir die \gesperrt{Differentiale} von $\alpha, \beta$. Die Differentiale der Konstanten und nur diese sind gleich Null. Die anderen Differentiale sind alle vollständig bestimmt, wenn eines von ihnen willkürlich (z. B. konstant) angenommen ist.

3. \gesperrt{Sind $\alpha, \beta, \gamma, \ldots$, beliebige Funktionen in $\Omega$, die einer Gleichung}
\weq{8}{F(\alpha, \beta, \gamma, \ldots) = 0}
\gesperrt{genügen, so ist}
\weq{9}{F'(\alpha)\,d\alpha + F'(\beta)\,d\beta + F'(\gamma)\,d\gamma + \cdots = 0.}

Um diesen Satz zu beweisen, nehme man einen Punkt $\mathfrak{P}$, in dem keine der Funktionen $\alpha, \beta, \gamma \ldots$ oder $d\alpha, d\beta, d\gamma, \ldots$ unendlich oder Null wird und auch die abgeleiteten $F'(\alpha), F'(\beta), F'(\gamma)$ nicht verschwinden, und ordnen wie in (3) die Funktion $F$ nach Potenzen und Produkten von $\alpha - \alpha_0, \beta - \beta_0, \gamma - \gamma_0, \ldots$ Dann ergibt sich, daß die Gleichung (9) für unendlich viele Punkte, und mithin identisch erfüllt ist. Als Spezialfälle von (9) ergeben sich:
% [sic] printed ``ordnen'' continuing ``nehme man'', a plural verb without subject; the impersonal construction requires ``ordne''
\weq{10}{\begin{aligned}
d(\alpha+\beta) &= d\alpha + d\beta,\\
d(\alpha\beta) &= \beta\,d\alpha + \alpha\,d\beta,\\
d\left(\frac{\alpha}{\beta}\right) &= \frac{\beta\,d\alpha - \alpha\,d\beta}{\beta^2}.
\end{aligned}}

Der Begriff des Differentials läßt sich ohne weiteres auf \gesperrt{holomorphe Funktionale} übertragen, wenn die Funktionalvariablen als Konstanten betrachtet werden. Der Satz 3. und die daraus folgenden Formeln (10) gelten dann auch noch, wenn $\alpha, \beta, \gamma, \ldots$ holomorphe Funktionale sind und $F(\alpha, \beta, \gamma, \ldots)$ eine ganze rationale Funktion von $\alpha, \beta, \gamma, \ldots$, ist.
